\noindent 本章按老师课件的 ERD 表达方式，将一组业务概念放在一张可读的关系图内。六组图与第 4 章六组类图逐组对应，基数、键和约束可与附录 C 复核。跨组引用以稳定 ID 连接，图中“目标”不代表本地数据库已创建。
% 007 数据设计：分模块 ERD；图内“目标”表不得误读为当前已落库。
% ER-CUST 顾客/预约。字段约束详见 usecase/数据字典.md。
\newcommand{\BSRERCust}{\begin{tikzpicture}[x=1cm,y=1cm]
  \BSRERCard{customers}{0}{0}{customers / users}{PK customer\_id\newline UQ email, NN name\newline 现有 users 映射}
  \BSRERCard{reservations}{5.2}{0}{reservations}{PK id\newline FK user\_id,store\_id\newline NN pickup\_at,return\_at}
  \BSRERCard{items}{10.4}{0}{reservation\_items}{PK id\newline FK reservation\_id,variant\_id\newline CK qty>0; NOK 快照}
  \BSRERCard{consents}{5.2}{-3.8}{damage\_consents}{PK/FK reservation\_id\newline FK user\_id\newline NN version,accepted\_at}
  \BSRERCard{requests}{0}{-3.8}{reservation\_requests}{PK actor/channel/key\newline UQ/FK reservation\_id\newline NN payload\_hash,created\_at}
  \BSRERCard{variant}{10.4}{-3.8}{variants}{PK id\newline style,size,skill\newline 规格而非物理单件}
  \draw[Navy] (customers) -- node[above,font=\tiny,fill=white]{1:n} (reservations);
  \draw[Navy] (reservations) -- node[above,font=\tiny,fill=white]{1:n} (items);
  \draw[Navy] (reservations) -- node[right,font=\tiny,fill=white]{1:0..1 同意} (consents);
  \draw[Navy] (requests.north east) -- node[above,sloped,font=\tiny,fill=white]{1:1} (reservations.south west);
  \draw[Navy,dashed] (customers) -- node[left,font=\tiny,fill=white]{actor FK} (requests);
  \draw[Navy] (items) -- node[right,font=\tiny,fill=white]{n:1 FK variant\_id} (variant);
  \node[bsrNote] at (5.2,-7.1) {租期是预约值字段；购物车不占库存。提交映射按实际账号、渠道和键唯一；历史订单可无映射。};
\end{tikzpicture}}

% ER-INV 商品/单件/门店；类别树是目标表。
\newcommand{\BSRERInventory}{\begin{tikzpicture}[x=1cm,y=1cm]
  \BSRERCard{category}{0}{0}{categories 目标}{PK id\newline FK parent\_id→本表\newline CK 无环}
  \BSRERCard{variant}{5.2}{0}{variants}{PK id\newline FK category\_id 目标\newline NN style,size; retail\_nok}
  \BSRERCard{unit}{10.4}{0}{units}{PK id; UQ inventory\_no\newline FK variant\_id,store\_id\newline FK home\_store\_id; status}
  \BSRERCard{store}{0}{-4}{stores}{PK id\newline NN name,address\newline open\_hours}
  \BSRERCard{policy}{5.2}{-4}{price\_policies}{PK id\newline FK variant\_id,store\_id\newline UQ variant/store/version}
  \BSRERCard{events}{10.4}{-4}{stock\_events}{PK id; FK unit\_id\newline from/to store,status\newline NN created\_at}
  \draw[Navy] (category) -- node[above,font=\tiny,fill=white]{1:n 类别} (variant);
  \draw[Navy] (variant) -- node[above,font=\tiny,fill=white]{1:n 单件} (unit);
  \draw[Navy] (category.west) to[out=165,in=115,looseness=2] node[left,font=\tiny]{自反 1:n} (category.north);
  \draw[Navy] (store) -- node[above,font=\tiny,fill=white]{1:n 价目} (policy);
  \draw[Navy] (policy) -- node[right,font=\tiny,fill=white]{n:1 规格} (variant);
  \draw[Navy] (unit) -- node[right,font=\tiny,fill=white]{1:n 事件} (events);
  \draw[Navy,dashed] (store) -- node[sloped,above,font=\tiny,fill=white]{当前/归属 两条 FK} (unit);
\end{tikzpicture}}

% ER-CONT 合同版本、预约行和单件 m:n 连接。
\newcommand{\BSRERContract}{\begin{tikzpicture}[x=1cm,y=1cm]
  \BSRERCard{reservation}{0}{0}{reservations}{PK id\newline FK user\_id,store\_id\newline status,contract\_version}
  \BSRERCard{contract}{5.2}{0}{contracts}{PK id; UQ code\newline FK reservation\_id\newline UQ reservation/version}
  \BSRERCard{clines}{10.4}{0}{contract\_lines 目标}{PK id; FK contract\_id\newline variant,qty,NOK\newline 当前在 snapshot\_json}
  \BSRERCard{prep}{0}{-3.8}{preparations}{PK/FK reservation\_id\newline assigned\_to,status\newline version}
  \BSRERCard{punit}{5.2}{-3.8}{preparation\_units}{PK (res,unit)\newline FK item,unit\newline UQ active unit}
  \BSRERCard{unit}{10.4}{-3.8}{units}{PK id; UQ inventory\_no\newline 状态/当前店/归属店}
  \BSRERCard{item}{0}{-7.6}{reservation\_items}{PK id\newline FK reservation,variant\newline qty,NOK 快照}
  \BSRERCard{runit}{5.2}{-7.6}{rental\_units}{PK (item,unit)\newline out\_at,returned\_at\newline FK return\_store\_id}
  \BSRERCard{adjust}{10.4}{-7.6}{contract\_adjustments}{PK id; FK reservation\_id\newline UQ reservation/reference\newline base\_version,差额 NOK}
  \draw[Navy] (reservation) -- node[above,font=\tiny,fill=white]{1:n 版本} (contract);
  \draw[Navy] (contract) -- node[above,font=\tiny,fill=white]{1:n 目标行} (clines);
  \draw[Navy] (reservation) -- node[left,font=\tiny,fill=white]{1:1 当前备货} (prep);
  \draw[Navy] (prep) -- node[above,font=\tiny,fill=white]{1:n} (punit);
  \draw[Navy] (punit) -- node[above,font=\tiny,fill=white]{n:1 实物} (unit);
  \draw[Navy] (punit) -- node[right,font=\tiny,fill=white]{交付后 1:1} (runit);
  \draw[Navy] (item) -- node[above,font=\tiny,fill=white]{1:n 逐件借还} (runit);
  \draw[Navy] (item) -- node[sloped,above,font=\tiny,fill=white]{1:n 准备绑定} (punit);
  % 调整记录跨过两条逐件连接线；以卡片 FK 和数据字典表达，避免连线压住字段。
\end{tikzpicture}}

% ER-SERV 归还事实、鉴定证据、维修和安全验收。
\newcommand{\BSRERService}{\begin{tikzpicture}[x=1cm,y=1cm]
  \BSRERCard{inspection}{0}{0}{inspections}{PK id\newline FK reservation\_id,unit\_id\newline UQ reservation+unit}
  \BSRERCard{case}{5.2}{0}{damage\_cases}{PK id; UQ/FK inspection\_id\newline FK assessor\_id\newline grade,status,version}
  \BSRERCard{photo}{10.4}{0}{damage\_photos}{PK id; FK case\_id\newline part,image BLOB NN\newline FK uploaded\_by}
  \BSRERCard{repair}{0}{-4}{maintenance}{PK id; FK unit\_id\newline FK contractor\_id\newline rate\_id,rate 快照}
  \BSRERCard{rates}{5.2}{-4}{contractor\_rates}{PK id; FK contractor\_id\newline effective\_from,hourly\_nok\newline UQ contractor/date}
  \BSRERCard{accept}{10.4}{-4}{safety\_acceptances 目标}{PK id; FK maintenance\_id\newline FK reviewer\_id\newline decision,reviewed\_at}
  \draw[Navy] (inspection) -- node[above,font=\tiny,fill=white]{1:0..1} (case);
  \draw[Navy] (case) -- node[above,font=\tiny,fill=white]{1:n 证据} (photo);
  \draw[Navy] (case) -- node[left,font=\tiny,fill=white]{0:1 工单} (repair);
  \draw[Navy] (rates) -- node[above,font=\tiny,fill=white]{1:n} (repair);
  \draw[Navy] (repair.south) -- ++(0,-.55) -| node[pos=.76,below,font=\tiny,fill=white]{1:n 验收轮次} (accept.south);
  \node[bsrNote] at (5.2,-7.2) {当前 `maintenance\_events` 保存验收事实；目标表须保留轮次。照片按案件对象授权，不能放公开目录。};
\end{tikzpicture}}

% ER-FIN 押金、超时、损坏和资金/补差流水。
\newcommand{\BSRERFinance}{\begin{tikzpicture}[x=1cm,y=1cm]
  \BSRERCard{runit}{0}{0}{rental\_units}{PK item\_id+unit\_id\newline late\_minutes/fee\_nok\newline 逐件实还冻结}
  \BSRERCard{case}{5.2}{0}{damage\_cases}{PK id\newline grade,fee\_nok\newline manager approved\_by}
  \BSRERCard{settle}{10.4}{0}{damage\_settlements}{PK/FK reservation\_id\newline damage/late/total NOK\newline deduct/refund/due NOK}
  \BSRERCard{payments}{0}{-4}{payments}{PK id; FK reservation\_id\newline type,amount\_nok,status\newline UQ idempotency\_key}
  \BSRERCard{contacts}{5.2}{-4}{damage\_contacts}{PK id; FK reservation\_id\newline FK actor\_id\newline channel,outcome,time}
  \BSRERCard{receipts}{10.4}{-4}{damage\_receipts}{PK id; FK reservation\_id\newline amount\_nok,reference\newline UQ reservation/reference}
  \draw[Navy] (runit.north) -- ++(0,.55) -| node[pos=.75,above,font=\tiny,fill=white]{n:1 超时} (settle.north);
  \draw[Navy] (case) -- node[above,font=\tiny,fill=white]{n:1} (settle);
  \draw[Navy] (settle) -- node[sloped,above,font=\tiny,fill=white]{1:n 资金事件} (payments);
  \draw[Navy] (settle) -- node[sloped,above,font=\tiny,fill=white]{1:n 沟通} (contacts);
  \draw[Navy] (settle) -- node[right,font=\tiny,fill=white]{1:n 补差凭证} (receipts);
  \node[bsrNote] at (5.2,-7.2) {费用=逐件超时+核定损坏；抵扣不超过押金；争议或回执不明不得标已到账。};
\end{tikzpicture}}

% ER-STAFF 角色单表、档案扩展、排班/打卡、提醒与审计。
\newcommand{\BSRERWorkforce}{\begin{tikzpicture}[x=1cm,y=1cm]
  \BSRERCard{users}{0}{0}{users}{PK id; UQ email\newline CK role 枚举\newline FK store\_id}
  \BSRERCard{profiles}{5.2}{0}{employee\_profiles}{PK/FK user\_id\newline UQ employee\_no\newline job\_title,qualification}
  \BSRERCard{shifts}{10.4}{0}{shifts}{PK id; FK user/store\_id\newline work\_date,start/end\newline UQ user+date}
  \BSRERCard{punch}{0}{-4}{attendance}{PK id; FK user\_id\newline clock\_in/clock\_out\newline UQ user+work\_date}
  \BSRERCard{alerts}{5.2}{-4}{order\_alert\_reads}{PK user\_id+reservation\_id\newline read\_at NN\newline 未解决看业务状态}
  \BSRERCard{audit}{10.4}{-4}{audit\_log}{PK id; FK actor\_id\newline action,target\_type/id\newline detail\_json,created\_at}
  \draw[Navy] (users) -- node[above,font=\tiny,fill=white]{1:0..1 扩展} (profiles);
  \draw[Navy] (users.north) -- ++(0,.55) -| node[pos=.72,above,font=\tiny,fill=white]{1:n 计划} (shifts.north);
  \draw[Navy] (users) -- node[left,font=\tiny,fill=white]{1:n 实际} (punch);
  \draw[Navy] (users) -- node[sloped,above,font=\tiny,fill=white]{1:n 本人已读} (alerts);
  \draw[Navy] (users) -- node[sloped,above,font=\tiny,fill=white]{1:n 操作} (audit);
  \node[bsrNote] at (5.2,-7.2) {当前单表 role；目标子类表需迁移。打卡不等于排班，已读提醒不等于订单闭环。};
\end{tikzpicture}}

\noindent 六图的共同约定：PK 为恒定主键，FK 为引用键，UQ 为唯一，NN 为非空，CK 为检查；目标表直接标“目标”。字段类型、迁移状态及规则来源见《数据字典》和《现有目标数据差异》。图仅展示主链路，省略重复的审计字段。
\BSRNewFigure{\BSRERCust}{ER-CUST 顾客、预约行与条款同意；UC-01～04}{fig:er-cust}
\BSRNewFigure{\BSRERInventory}{ER-INV 商品类别、规格、实体单件与门店；UC-02/07/09/13}{fig:er-inv}
\BSRNewFigure{\BSRERContract}{ER-CONT 合同版本、准备和逐件交接；UC-05～09}{fig:er-cont}
\BSRNewFigure{\BSRERService}{ER-SERV 归还鉴定、照片、维修与安全验收；UC-09/10/12}{fig:er-serv}
\BSRNewFigure{\BSRERFinance}{ER-FIN 超时与受损费用、押金、补差和资金回执；UC-03/08/11}{fig:er-fin}
\BSRNewFigure{\BSRERWorkforce}{ER-STAFF 员工档案、排班、考勤、提醒与审计；UC-14/15}{fig:er-staff}
\begin{BSRTable}
\rowcolor{TableHead}\color{white}\bfseries 关系或约束 & \color{white}\bfseries 设计理由与实现边界 \\
顾客与预约 / 预约与行 & 均为 1:n，FK 放预约和行；条款同意保留版本，历史事件不可覆盖。\\
类别树 / 账号档案 & 类别 parent\_id 是自反 FK，并通过服务规则阻止环；员工档案以 user\_id 同时作 PK/FK，顾客账号不强制建档。\\
预约行与实体装备 & 准备绑定和实际租用分别用连接表拆解 m:n，保存时间与责任人；单件活跃占用有唯一/排他守卫。\\
费用与签署快照 & 参考价值 8\% 押金、价目版本、合同版本、技师费率和规则版本在交易时冻结；后改价不回写历史单。\\
跨店归还与调拨 & 实还门店记录在逐件交接；单件当前店可改变，归属店只随正式调拨接收更新。\\
目标与现有 & categories、contract\_lines、safety\_acceptances 等是目标关系；当前演示仍使用文本类别、JSON 快照或事件表。\\
\end{BSRTable}

\subsection{6.1 两客户端身份与报价关系补图}
\noindent 本图只补充渠道身份和短时报价关系；顾客、门店、规格、物理单件及预约仍以本章六组 ERD 为准。活跃绑定采用部分唯一索引，历史解绑行仍留存；购物车版本与报价指纹用于防止另一端更新后误下旧单。
\BSRDiagram{019-erd-identity}{ER-MINI-A 微信身份、一次性凭证与独立会话；PK/FK 及活跃唯一约束}{fig:er-mini-identity}
\BSRDiagram{019-erd-commerce}{ER-MINI-B 共用顾客、购物车、报价与预约；版本及归属约束}{fig:er-mini-commerce}
\begin{BSRTable}
\rowcolor{TableHead}\color{white}\bfseries 关系 / 约束 & \color{white}\bfseries 可检验的物理设计 \\
顾客与绑定 & wechat\_bindings.user\_id 外键指向 users.id；同一 AppID 的活跃 OpenID 摘要唯一、活跃顾客唯一。解绑只填时间，不删除历史。\\
绑定与小程序会话 & 1:n；mini\_sessions.binding\_id、user\_id 指向活跃关系。令牌仅保存摘要；账号认证纪元或绑定状态变化后拒绝旧令牌。\\
顾客与购物车/报价 & customer\_carts.customer\_id 唯一，购物车行以 (cart\_id,variant\_id) 为复合主键；mini\_quotes.customer\_id 关联顾客，报价短期有效，指纹绑定完整输入。\\
购物车与预约 & 购物车不是订单外键；经核价与库存事务后才创建 reservations.user\_id 所属的预约。历史订单保存独立政策/金额快照，购物车后续变化不回写。\\
一次性凭证 & mini\_login\_challenges 和 mini\_code\_uses 用不可逆摘要与过期/消耗字段；不把临时码或平台会话密钥当持久业务键。\\
\end{BSRTable}

\subsection*{交付边界与待确认事项}
\addcontentsline{toc}{subsection}{交付边界与待确认事项}
\begin{BSRTable}
\rowcolor{TableHead}\color{white}\bfseries 项目 & \color{white}\bfseries 当前状态 \\
现状证据 & 案例明确 Excel 工作簿、纸质维修记录等；不足以证明的交付、结算、调拨顺序保留为证据缺口。消息线型不表示证据等级；不以补画箭头代替访谈。\\
商业政策 & 8\% 参考价值押金、首单 2 小时宽限及等级扣费是本项目演示政策；正式费率、取消、赔偿和退款条款待 BSR 批准与版本化。\\
外部接口 & 真实支付、银行退款、电子签署、汇率及通知服务为目标集成；本地演示不能作为资金到账证据。定位变更只更新参考展示币种，合同以挪威市场 NOK 冻结。\\
生产验证 & 并发、容量、安全、灾备与第三方失败恢复尚需专门测试和组织验收；本地自动测试只证明对应演示路径。\\
\end{BSRTable}
\subsection*{参考资料与版本依据}
\addcontentsline{toc}{subsection}{参考资料与版本依据}
\noindent BSR 案例、课程课件与课程总结提供业务背景和建模方法；初版设计需求文档是历史版本；QuickBooks 与 Shopify ERD 是图面组织参考；参考论文只提供版式参考。公开系统用于功能与交互对照，不替代 BSR 组织事实。完整资料清单见附录 F；旧新用例语义对照见附录 E。
